% Diagram of the proof of Theorem B (Part 3). Boxes are results; an arrow X -> Y means that the proof of Y uses X.
\begin{figure}[tbp]
\centering
\begin{tikzpicture}[x=1cm,y=1cm]
% top: Theorem B and the spanning argument
\node[pfmain] (B) at (0,0) {Thms.~\ref{main:thm:B}, \ref{sp-moment:thm:B}\\Theorem~B};
\node[pfnode, text width=2.0cm] (RAM) at (-4.4,-0.6) {Lem.~\ref{ub-divis:lem:ram}\\ramified forms (Part~2)};
\node[pfnode] (SP) at (0,-1.8) {Props.~\ref{sp-moment:prop:one}, \ref{sp-moment:prop:few}, Thm.~\ref{sp-moment:thm:spread}\\spanning argument};
% row 2: the inputs of the spanning argument
\node[pfnode] (PET) at (-6.3,-3.7) {Thms.~\ref{sp-unfold:thm:series}, \ref{sp-unfold:thm:petersson}\\Petersson products};
\node[pfnode, text width=2.0cm] (DIM) at (-1.275,-3.7) {Prop.~\ref{sp-moment:prop:dim}\\dimension of the target};
\node[pfnode, text width=2.0cm] (GOOD) at (1.275,-3.7) {Thm.~\ref{sp-unfold:thm:good}\\good forms};
\node[pfnode, text width=2.0cm] (MOM) at (3.825,-3.7) {Thm.~\ref{sp-unfold:thm:moment}\\second moment};
\node[pfnode, text width=2.0cm] (BAD) at (6.375,-3.7) {Thm.~\ref{sp-main:thm:bad}\\bad share};
% row 3
\node[pffact, text width=2.0cm] (FUN) at (-6.375,-5.6) {Fact~\ref{sp-unfold:fact:unfold}\\unfolding};
\node[pfnode, text width=2.0cm] (SYS) at (-3.825,-5.6) {Thm.~\ref{sp-moment:thm:system}\\Bessel systems};
\node[pfnode, text width=2.0cm] (MULT) at (-1.275,-5.6) {Thm.~\ref{sp-moment:thm:trace}, Cor.~\ref{sp-moment:cor:mult}\\multiplicities};
\node[pffact, text width=2.0cm] (FDIM) at (1.275,-5.6) {Fact~\ref{sp-moment:fact:dim}\\algebraic forms};
\node[pfnode, text width=2.0cm] (MAIN) at (3.825,-5.6) {Thm.~\ref{sp-main:thm:main}\\main term};
\node[pfnode, text width=2.0cm] (REM) at (6.375,-5.6) {Thm.~\ref{sp-remain:thm:main}\\remaining terms};
% row 4
\node[pfnode, text width=2.0cm] (REAL) at (-6.375,-7.5) {Thm.~\ref{sp-frames:thm:real}\\local bounds at~$\mathbb{C}$};
\node[pfnode, text width=2.0cm] (STAT) at (-3.825,-7.5) {Thms.~\ref{sp-moment:thm:laplace}, \ref{sp-moment:thm:stationary}\\Laplace's method,\\stationary point};
\node[pffact, text width=2.0cm] (FTR) at (-1.275,-7.5) {Fact~\ref{sp-moment:fact:trace}\\comparison of traces};
\node[pfnode, text width=2.0cm] (SIEVE) at (1.275,-7.5) {Thm.~\ref{sp-main:thm:sieve}\\the sieve};
\node[pffact, text width=2.0cm] (FSIEG) at (3.825,-7.5) {Fact~\ref{sp-main:fact:siegel}\\Siegel's formula};
\node[pfnode, text width=2.0cm] (FIN) at (6.375,-7.5) {\S\ref{sp-frames:ss:finite}\\local estimates at finite places};
% row 5
\node[pffact, text width=2.0cm] (FPOI) at (1.275,-9.2) {Fact~\ref{sp-main:fact:poisson}\\Poisson summation};
% arrows: premise -> conclusion
\draw[pfedge] (RAM) -- (B);
\draw[pfedge] (SP) -- (B);
\draw[pfedge] (PET) |- (SP);
\draw[pfedge] (DIM) -- (SP);
\draw[pfedge] (GOOD) -- (SP);
\draw[pfedge] (MOM) -- (SP);
\draw[pfedge] (BAD) |- (SP);
\draw[pfedge] (FDIM) -- (DIM);
\draw[pfedge] (MULT) -- (DIM);
\draw[pfedge] (SYS) -- (DIM);
\draw[pfedge] (SYS) -- (PET);
\draw[pfedge] (FUN) -- (PET);
\draw[pfedge] (FTR) -- (MULT);
\draw[pfedge] (REAL) -- (SYS);
\draw[pfedge] (STAT) -- (SYS);
\draw[pfedge] (MAIN) -- (MOM);
\draw[pfedge] (REM) -- (MOM);
\draw[pfedge] (SIEVE) -- (MAIN);
\draw[pfedge] (FSIEG) -- (MAIN);
\draw[pfedge] (FIN) -- (MAIN);
\draw[pfedge] (FPOI) -- (SIEVE);
\end{tikzpicture}
\caption{The proof of Theorem~B (Part~3, Sections~\ref{sec:sp-moment}--\ref{sec:sp-remain}). Each
box is a result, with its number and a short name; a box with several numbers, or with the number
of a subsection, stands for a group of results. An arrow from one box to another means that the proof of the second, as written in
this paper, uses the first. Dashed grey boxes are background results admitted without proof.}
\label{fig:part3}
\end{figure}
